\subsection{1D Sod Shock Tube}
\label{SEC:Testcases:Sod}

\begin{figure}[h!tb]
  \begin{centering}
    \begin{subfigure}{3.0in}
      \insertimage[width=\textwidth]{test_cases/sod_start}
      \caption{Initial condition.}
      \label{FIG:SodStart}
    \end{subfigure}%
    \begin{subfigure}{3.0in}
      \insertimage[width=\textwidth]{test_cases/sod_final}
      \caption{\PO[3] result at time 0.1712 seconds.}
      \label{FIG:SodFinal}
    \end{subfigure}
    \caption{Sod shock tube test case.}
    \label{FIG:TestcaseSod}
  \end{centering}
\end{figure}

The Sod shock tube is a simple case for evaluating shock capturing.
%
When using the internally generated grid, the domain will be 1.0 meter
long (centered at the origin) and the left\slash right boundary conditions
will be set to extrapolation.
%
For the initial condition at time $t$=0, imagine that there is an
invisible diaphram in the center of the tube as shown in \fref{FIG:SodStart}.
%
The flow on the left side is set to 1 \unitPres{} and 1 \unitDens{}, and the
flow on the right has 1/\nth{10} the pressure and 1/\nth{8} the density.
Neither side has any initial velocity.
 
Once the simulation begins, the diaphram is removed and the high-pressure fluid
on the left moves to the right.
%
This forms a shock wave moving to the right.
%
Somewhat behind the shock wave, there is a contact discontinuity in density that
also moves to the right.
%
Meanwhile an expansion fan forms on the left side and moves to the left.
%
\Fref{FIG:SodFinal} shows a sample \PO[3] result on a grid of 100 cells using
the resolution marker to identify troubled cells and applying a filter with
$\inp[nmlLimiter::LocalFilterOrder]{Local\_Filter\_Order}=4$ to damp
oscillations.

\SeeAlso{
\inp[nmlInitialization::Initialization]{Initialization}=23,
\inp[nmlGrid::InternalGridDoF]{Internal\_Grid\_DoF},
\inp[nmlTestcaseSod]{nml\_testcase\_sod}
}

